% Created 2026-09-07 Mon 06:26
% Intended LaTeX compiler: pdflatex
\documentclass[11pt, a4paper, abstract=yes]{scrartcl}
\usepackage[utf8]{inputenc}
\usepackage[T1]{fontenc}
\usepackage{graphicx}
\usepackage{longtable}
\usepackage{wrapfig}
\usepackage{rotating}
\usepackage[normalem]{ulem}
\usepackage{amsmath}
\usepackage{amssymb}
\usepackage{hyperref}
\usepackage{amsthm,amssymb,amsmath}
\usepackage{hyperref}
\usepackage{graphicx}
\theoremstyle{definition}
\newtheorem{definition}{Definition}
\newtheorem{mexample}[definition]{Example}
\theoremstyle{remark}
\newtheorem{remark}[definition]{Remark}
\theoremstyle{plain}
\newtheorem{lemma}[definition]{Lemma}
\newtheorem{claim}[definition]{Claim}
\newtheorem{proposition}[definition]{Proposition}
\newtheorem{theorem}[definition]{Theorem}
\newtheorem{corollary}[definition]{Corollary}
\newtheorem{cor}[definition]{Corollary}
\newtheorem{conjecture}[definition]{Conjecture}
\newtheorem{question}[definition]{Question}
\newcommand{\set}[1]{\left\{{#1}\right\}}
\newcommand\setsuchas[2]{\left\{\,{#1}\,\vrule\,{#2}\,\right\}}
\newcommand{\vek}[1]{\boldsymbol{#1}}
\newcommand{\NN}{{\mathbb{N}}}
\newcommand{\NNp}{{\mathbb{N}_+}}
\newcommand{\ZZ}{{\mathbb{Z}}}
\newcommand{\CC}{{\mathbb{C}}}
\newcommand{\RR}{{\mathbb{R}}}
\newcommand{\QQ}{{\mathbb{Q}}}
\newcommand{\seqnum}[1]{\href{https://oeis.org/#1}{\textrm \underline{#1}}}
\author{Jan Snellman\thanks{Department of Mathematics, Linköping University, 581 83 Linköping, Sweden. \texttt{jan.snellman@liu.se}. ORCID: \href{https://orcid.org/0009-0002-6676-5068}{0009-0002-6676-5068}}}
\usepackage{iftex}
\ifPDFTeX\usepackage[T1]{fontenc}\else\usepackage{fontspec}\fi
\usepackage{orcidlink}
\hypersetup{hidelinks}
\newcommand{\Koz}{\operatorname{Koz}}
\newcommand{\shift}{\operatorname{shift}}
\newcommand{\ff}{\operatorname{ff}}
\newcommand{\IM}{\operatorname{IM}}
\newcommand{\inn}{\operatorname{in}}
\newcommand{\gin}{\operatorname{gin}}
\newcommand{\conv}{\operatorname{conv}}
\newcommand{\indeg}{\operatorname{indeg}}
\author{Jan Snellman \orcidlink{0009-0002-6676-5068}}
\date{\today}
\title{Hilbert functions over the exterior algebra and \(f\)-vectors in higher rank\\\medskip
\large III: Lattice points of the Kozlov polytope, and a corrected Kruskal--Katona theorem}
\hypersetup{
 pdfauthor={Jan Snellman \orcidlink{0009-0002-6676-5068}},
 pdftitle={Hilbert functions over the exterior algebra and \(f\)-vectors in higher rank},
 pdfkeywords={},
 pdfsubject={},
 pdfcreator={Emacs 30.2 (Org mode 9.7.11)}, 
 pdflang={English}}
\begin{document}

\maketitle
\begin{abstract}
{[}TO WRITE.] Companion to Part I of this series, \emph{Minkowski sums of Kozlov simplices}, which determines the convex hull \(\Koz(n;d_1,\dots,d_r)\) of
the Hilbert functions of quotients of a graded free module over the exterior algebra by a proper
monomial submodule. That result is a statement about a polytope; it says nothing about which of the
polytope's lattice points are actually achievable. Here we answer that finer question. The
characterization is a rank-\(r\) growth bound generalizing Kruskal--Katona --- fill the generators of
largest degree to capacity first, then apply the shadow operator to the remainder --- and the
extremal representative of each achievable Hilbert function is the unique lex submodule realizing
it. This is Amata and Crupi's theorem, which we state in a \textbf{repaired} form: the published version is
false as printed whenever two or more free generators tie at the initial critical degree, and we give
the smallest counterexample, the correct clause, and an exhaustive verification.
\end{abstract}

\setcounter{tocdepth}{2}
\tableofcontents
\section{The vertex and facet structure of the Kozlov polytope}
\label{sec:org6a57ba1}
The lattice-point question of this paper is posed inside the polytope \(\Koz(n;\vec d)\) of Part I,
and answering it needs that polytope described precisely rather than merely bounded. This section
collects what is known: the vertex description at every rank for the all-ones right-angle simplex,
and the facet description at every rank. Both were developed while Part I was being written and are
reproduced here because this is where they are used.
\subsection{The pairwise reduction, proved for the all-ones simplex}
\label{sec:orga5ed02d}
The extremal tensor of Definition \texttt{def-angle} (Part I) is a rank-\(r\) object, but for \(a=(1,\dots,1)\) it
turns out to be determined by its rank-two restrictions, and indeed by the consecutive ones alone.
This subsection proves that. Throughout it the leg vector is \(a=(1,\dots,1)\) and the shifts are
sorted, so that \(d_1\le\dots\le d_r\). Define the \emph{window} of summand \(p\) to be the integer
interval
\[
  W_p=[\,d_p+1,\;d_p+n\,],
\]
of length exactly \(n\). Given a choice \(\vec\jmath=(j_1,\dots,j_r)\), write \(m_p=d_p+j_p\), so
that \(m_p\in W_p\), and let \(D=D(\vec\jmath)\) be the digraph on \(W_1\cup\dots\cup W_r\) whose
edges are \(m_p\to x\), one for every index \(p\) and every \(x\in W_p\setminus\{m_p\}\).

The first lemma trades the polytope for that digraph.

\begin{lemma}
Let \(a=(1,\dots,1)\). Then \(M[\vec\jmath]=1\) if and only if \(D(\vec\jmath)\) is acyclic.
\label{orgf5c2e7b}
\end{lemma}

\begin{proof}
Let \(\varphi\) be a linear functional. Summing coordinates gives
\[
  \bigl\langle\varphi,\ \shift^{d_p}(v_j)\bigr\rangle=\sum_{i=1}^{j}\varphi_{d_p+i}
  =\Phi(d_p+j)-\Phi(d_p),
  \qquad
  \Phi(m):=\sum_{i\le m}\varphi_i .
\]
The offset \(\Phi(d_p)\) does not depend on \(j\), so \(\varphi\) exposes the vertex \(v_{j_p}\) of
summand \(p\) uniquely precisely when \(\Phi\) attains a strict maximum over \(W_p\) at \(m_p\). By
the normal-fan criterion, \(M[\vec\jmath]=1\) exactly when some single \(\varphi\) does this for
every \(p\) at once. Passing from \(\varphi\) to \(\Phi\) is a bijection onto real sequences, since
\(\varphi\) is recovered by differencing, and the coordinates of \(\Phi\) outside
\(W_1\cup\dots\cup W_r\) are unconstrained. Hence the condition is the solvability of the system
\[
  \Phi(m_p)>\Phi(x)
  \qquad\text{for every }p\text{ and every }x\in W_p\setminus\{m_p\},
\]
whose constraints are exactly the edges of \(D\).

It remains to observe that a system of strict inequalities \(X(u)>X(v)\), one for each edge of a
finite digraph, is solvable if and only if that digraph is acyclic. A cycle
\(u_0\to u_1\to\dots\to u_{k-1}\to u_0\) would force \(X(u_0)>X(u_1)>\dots>X(u_0)\), which is
impossible. Conversely an acyclic digraph admits a topological order, and setting \(X(w)\) to be
minus the position of \(w\) in that order satisfies every constraint.
\end{proof}

The next lemma is the one place where the windows being intervals of \emph{equal} length is used.

\begin{lemma}
For indices \(p\le s\le q\) one has \(W_p\cap W_q\subseteq W_s\).
\label{org484b196}
\end{lemma}

\begin{proof}
Assume the intersection is nonempty, the inclusion being trivial otherwise, and let
\(x\in W_p\cap W_q\). From \(x\in W_q\) and \(d_s\le d_q\) we get \(x\ge d_q+1\ge d_s+1\), and from
\(x\in W_p\) and \(d_p\le d_s\) we get \(x\le d_p+n\le d_s+n\). Hence \(x\in W_s\).
\end{proof}

\begin{lemma}
If \(D\) has a cycle, then \(D\) has a cycle of length two; that is, there are indices \(p\ne q\)
with \(m_q\in W_p\), \(m_p\in W_q\) and \(m_p\ne m_q\).
\label{org0e7b129}
\end{lemma}

\begin{proof}
Every edge \(m_p\to x\) has \(x\ne m_p\), so \(D\) has no loops and every cycle has length at least
two. Choose a cycle of minimum length \(k\), and suppose \(k\ge3\). Write it as
\(v_0\to v_1\to\dots\to v_{k-1}\to v_0\) with indices read modulo \(k\); a cycle of minimum length
is vertex-simple, so the values \(v_0,\dots,v_{k-1}\) are pairwise distinct. Each edge
\(v_i\to v_{i+1}\) is witnessed by at least one summand, and we fix one, calling it \(p(i)\), so
that \(m_{p(i)}=v_i\) and \(v_{i+1}\in W_{p(i)}\).

We claim first that \(W_{p(i)}\) contains no cycle value other than \(v_i\) and \(v_{i+1}\). Suppose
\(v_j\in W_{p(i)}\) for some \(j\not\equiv i,i+1\). The values are distinct, so
\(v_j\ne v_i=m_{p(i)}\) and therefore \(v_i\to v_j\) is an edge. Following that edge by the arc of
the original cycle running from \(v_j\) back to \(v_i\) produces a closed walk of length
\(1+((i-j)\bmod k)\). Because \(j\not\equiv i+1\) we have \((i-j)\bmod k\le k-2\), so the walk has
length at most \(k-1\); as every closed walk contains a directed cycle no longer than itself, this
contradicts the minimality of \(k\).

Now \(W_{p(i)}\) is an integer interval containing both \(v_i\) and \(v_{i+1}\), hence containing
every integer between them, and by the previous paragraph it contains no further cycle value. So no
cycle value lies strictly between \(v_i\) and \(v_{i+1}\), which is to say that \(v_i\) and
\(v_{i+1}\) are neighbours in the increasing order of the \(k\) distinct cycle values. Applying this
at the largest value \(v_{i^\ast}\), both of its cycle-neighbours \(v_{i^\ast-1}\) and
\(v_{i^\ast+1}\) are order-adjacent to the maximum and smaller than it, so each is the
second-largest value. Since \(k\ge3\) the indices \(i^\ast-1\) and \(i^\ast+1\) are distinct modulo
\(k\), making \(v_{i^\ast-1}\ne v_{i^\ast+1}\); two distinct values cannot both be second-largest.

Therefore \(k=2\), and the cycle reads \(v\to w\to v\) with \(v\ne w\). Some \(p\) witnesses the
first edge, so \(m_p=v\) and \(w\in W_p\), and some \(q\) witnesses the second, so \(m_q=w\) and
\(v\in W_q\). Finally \(m_p=v\ne w=m_q\) forces \(p\ne q\).
\end{proof}

Only the interval property of the windows entered that argument. Neither the sorting of the shifts
nor the equality of the window lengths was used, so Lemma \ref{org0e7b129} holds verbatim for an
arbitrary finite family of integer intervals with designated points. The next lemma is where those
hypotheses become indispensable.

\begin{lemma}
If \(D\) has a cycle of length two, then \(D\) has one at a consecutive pair \((p,p+1)\).
\label{orgc85a815}
\end{lemma}

\begin{proof}
We argue by strong induction on the gap \(g=q-p\) of a length-two cycle \((p,q)\) with \(p<q\), so
that \(m_q\in W_p\), \(m_p\in W_q\) and \(m_p\ne m_q\). There is nothing to prove when \(g=1\), so
let \(g\ge2\) and fix an index \(s\) with \(p<s<q\). Both \(m_p\) and \(m_q\) lie in
\(W_p\cap W_q\), so Lemma \ref{org484b196} places both of them in \(W_s\). Three exhaustive cases remain.

Suppose first that \(m_s\in W_p\) and \(m_s\ne m_p\). Then the pair \((p,s)\) is a length-two cycle,
since \(m_s\in W_p\) by assumption and \(m_p\in W_s\) as just observed. Its gap \(s-p\) is smaller
than \(g\), so the induction applies.

Suppose next that \(m_s=m_p\). Then the pair \((s,q)\) is a length-two cycle: we have \(m_q\in W_s\)
as observed, while \(m_s=m_p\in W_q\) by hypothesis and \(m_s=m_p\ne m_q\). Its gap \(q-s\) is
smaller than \(g\).

Suppose finally that \(m_s\notin W_p\). From \(m_s\ge d_s+1\ge d_p+1\) and
\(W_p=[d_p+1,d_p+n]\) we get \(m_s>d_p+n\). On the other hand \(m_q\in W_p\) gives
\(m_q\le d_p+n\), so \(m_s>m_q\); in particular \(m_s\ne m_q\) and \(m_s>m_q\ge d_q+1\). Combining
this with \(m_s\le d_s+n\le d_q+n\) shows \(m_s\in W_q\), so the pair \((s,q)\) is again a
length-two cycle, of gap \(q-s<g\).
\end{proof}

\begin{theorem}
Let \(a=(1,\dots,1)\), let \(n\ge1\) and \(r\ge2\), let \(0\le d_1\le\dots\le d_r\), and let
\(\vec\jmath=(j_1,\dots,j_r)\) be any choice of vertices. The following are equivalent.
\begin{enumerate}
\item \(M[\vec\jmath]=1\).
\item \(M^{(p,q)}[j_p][j_q]=1\) for every pair \(p<q\).
\item \(M^{(p,p+1)}[j_p][j_{p+1}]=1\) for every \(p=1,\dots,r-1\).
\end{enumerate}

Here \(M^{(p,q)}\) denotes the rank-two extremal matrix of the summands \(p\) and \(q\) alone.
\label{org12ecee7}
\end{theorem}

\begin{proof}
Write \(D(p,q)\subseteq D\) for the digraph built from summands \(p\) and \(q\) alone. Lemma
\ref{orgf5c2e7b} applies at every rank, so \(M[\vec\jmath]=1\) if and only if \(D\) is acyclic, and
\(M^{(p,q)}[j_p][j_q]=1\) if and only if \(D(p,q)\) is acyclic.

For the implication from the first statement to the second, note that a cycle of \(D(p,q)\) is a
cycle of \(D\). The second statement implies the third because the third imposes fewer conditions.
For the remaining implication we argue by contraposition. If \(D\) has a cycle, then Lemma
\ref{org0e7b129} supplies a length-two cycle and Lemma \ref{orgc85a815} a consecutive one, say at
\((p,p+1)\). The two edges \(m_p\to m_{p+1}\) and \(m_{p+1}\to m_p\) both lie in \(D(p,p+1)\), which
is therefore not acyclic, so \(M^{(p,p+1)}[j_p][j_{p+1}]=0\) and the third statement fails.
\end{proof}

Theorem \ref{org12ecee7} is what the term \emph{pairwise reduction} will mean below. Its content is
that for the all-ones simplex the constraint network on the indices \(j_1,\dots,j_r\) is a \emph{path}
rather than a complete graph. Combined with the closed form of the rank-two matrix in Theorem
\texttt{thm-ones} (Part I), it determines the extremal tensor of \(\sum_i\shift^{d_i}(\angle(1,\dots,1))\) at every
rank, and hence the complete vertex description of that polytope for every \(r\) and every
\(\vec d\).

The scope of this deserves emphasis, since the corresponding statement for the Kozlov vector is
what Remark \texttt{obs-pairwise} (Part I) leaves open. Theorem \ref{org12ecee7} does \emph{not} extend Corollary
\texttt{cor-vdesc} (Part I), which concerns the Kozlov vector: for the Kozlov polytope \(\Koz(n;\vec d)\) the vertex
description is proved at rank two and remains conditional beyond it. Two remarks on the sharpness
of the all-ones result are worth recording. The hypothesis of equal
window lengths cannot be dropped from Lemma \ref{orgc85a815}: for the three intervals \([1,3]\),
\([1,4]\), \([2,3]\) with designated points \(2,4,3\) the only length-two cycle joins the first and
third, which are not consecutive. And the reduction genuinely needs the sorting, since the notion of
a consecutive pair depends on it.
\subsection{The facet description at general rank}
\label{sec:org970dea6}
Everything above describes vertices. The facets have so far been settled only at rank one, in
Corollary \texttt{cor-facets} (Part I). This subsection records what the same classical machinery gives for the
Minkowski sum at any rank, which is more than one might expect: the candidate facet normals are
determined by the summands, and the offsets are free.

\begin{proposition}
Let \(P_1,\dots,P_r\) be polytopes and \(P=P_1+\dots+P_r\). For every linear functional \(\varphi\),
the face of \(P\) on which \(\varphi\) is maximized satisfies
\[
  F(\varphi,P)=F(\varphi,P_1)+\dots+F(\varphi,P_r),
\]
the normal fan of \(P\) is the common refinement of the normal fans of the \(P_i\), and consequently
every facet normal of \(P\) lies on a wall of the normal fan of at least one summand.
\label{orgef6f203}
\end{proposition}

\begin{proof}
The displayed identity follows from additivity of support functions,
\[
  h(\varphi,P)=\sum_{i=1}^{r} h(\varphi,P_i),
\]
and it holds for any number of summands. Two functionals expose
the same face of \(P\) precisely when they expose the same face of every summand, which is the
statement that the normal fan of \(P\) refines each \(N(P_i)\) and is the coarsest fan doing so, that
is, their common refinement; refinement being associative, the case of general \(r\) follows from the
case \(r=2\). A maximal cone of the refinement is an intersection \(C_1\cap\dots\cap C_r\) of maximal
cones, and a wall of such an intersection is a wall of one of the \(C_i\). Since facets correspond to
walls of the normal fan, the last statement follows. This is Ziegler, \emph{Lectures on Polytopes}, Proposition 7.12, stated there for two summands; the general case follows by associativity of refinement, as above.
\end{proof}

Proposition \ref{orgef6f203} turns the search for an \(H\)-description of \(R\) into a finite problem,
but it must be applied with care, and one tempting shortcut is wrong. Facets of a polytope correspond
to \emph{rays} of its normal fan, and for \emph{full-dimensional} summands the rays of a common refinement are
the union of the summands' rays, so the facet normals of the sum would be the union of the summands'
facet normals --- the \(n\) inequalities of Corollary \texttt{cor-facets} (Part I) per summand, at most \(rn\) in all.
\textbf{That argument fails here.} Our summands are not full-dimensional, and, what matters, they have
\emph{different} lineality spaces; intersecting cones with different lineality spaces creates rays that
belong to no summand's fan. What Proposition \ref{orgef6f203} does still give is the weaker necessary
condition that a facet normal of \(R\) lies on a wall of some summand's fan, and walls correspond to
edges rather than facets, so the a priori bound is \(r\binom n2\).

The failure is not academic. Already for \(n=3\) and \(\vec d=(0,1)\), where \(R\) has six facets, only
four of the six inequalities of Corollary \texttt{cor-facets} (Part I) define facets of \(R\); the other two cut faces
of dimension one. The two missing facets are of a different shape, namely the bounds
\(x_{d_p+1}\ge a_1\) and a matching upper bound one coordinate past the end of the preceding window.
So the facet set of \(R\) is not the set of summand facets, and a correct list must be found rather
than inherited. For each candidate \(u\) the supporting value is For each candidate \(u\) the supporting value is
\[
  h(u,R)=\sum_{i=1}^{r} h\bigl(u,\shift^{d_i}(\angle(a))\bigr),
\]
a sum of maxima over known vertex sets. What remains is to decide which candidates are redundant.

One warning is needed, because it is a genuine trap rather than a formality. Each summand
\(\shift^{d_i}(\angle(a))\) has dimension \(n-1\) inside \(\mathbb R^N\), so its normal fan has a
large lineality space and a "facet normal of a summand" is a coset rather than a direction; a
particular inequality representation chooses a representative, and that representative need not be
the one cutting \(R\). The candidates must therefore be taken modulo each summand's affine hull. The
worked example below shows this happening.

\begin{mexample}
For \(n=3\), so \(a=(3,3,1)\), and \(\vec d=(0,1,2)\) in \(\mathbb R^5\), the polytope \(R\) has
dimension \(4\), sixteen vertices and exactly nine facets, one for each of the \(rn=9\) candidates:
none is redundant. The outward normals and their supporting values are
\[
  \begin{array}{lrrrr}
    u & h(u,R) & h_1 & h_2 & h_3\\[2pt]
    (0,1,0,0,0)   & 6   & 3 & 3 & 0\\
    (0,-1,3,0,0)  & 15  & 0 & 6 & 9\\
    (0,0,0,0,1)   & 1   & 0 & 0 & 1\\
    (0,0,0,1,0)   & 4   & 0 & 1 & 3\\
    (0,0,0,0,-1)  & 0   & 0 & 0 & 0\\
    (0,-1,0,0,0)  & -3  & 0 & -3& 0\\
    (0,0,-1,3,0)  & 6   & 0 & 0 & 6\\
    (0,0,-1,0,0)  & -3  & 0 & 0 & -3\\
    (0,0,0,-1,3)  & 0   & 0 & 0 & 0
  \end{array}
\]
The row \((0,0,0,0,-1)\) is the third summand's facet \(x_5\ge0\), inherited unchanged. The row
\((0,1,0,0,0)\), giving \(x_2\le6\), is the first summand's facet \(x_1-x_2\ge0\) read modulo its
affine hull \(x_1=3\), where it becomes \(x_2\le3\); the second summand lies in \(x_2=3\) and
contributes another \(3\). The computation was carried out independently in two systems with no
shared convex-hull library, and they agree on the dimension, both counts, every normal and every
offset.
\label{org0ea5801}
\end{mexample}

\begin{conjecture}
For \(a\) the Kozlov vector, let \(s=\lvert\lbrace d_1,\dots,d_r\rbrace\rvert\) be the number of
\emph{distinct} shifts. Then \(R=R(a;\vec d)\) of Definition \texttt{def-angle} (Part I) has exactly \(sn\) facets: no
candidate normal supplied by Proposition \ref{orgef6f203} is redundant. The same count holds for
\(\Koz(n;\vec d)\), which differs from \(R\) only by the translation and embedding of Remark
\texttt{rem-coordinates} (Part I).
\label{org36a14fc}
\end{conjecture}

The count is \(sn\) and not \(rn\) because tied shifts give \emph{identical} summands, whose walls
coincide and therefore supply the same candidates; this is the same collapsing that forces tied
shifts to carry equal vertex indices. At the extreme \(\vec d=(0,\dots,0)\) the sum is the dilation
\(r\angle(a)\), with \(n\) facets however large \(r\) is. The evidence is exact facet counts in
\(47\) configurations, with \(n=3,\dots,6\), ranks up to four, and \(\vec d\) ranging over patterns
that include ties \((0,1,1)\), \((0,0,1,1)\), \((0,2,2,4)\) and wide gaps \((0,2,4)\); there was no
exception.

For \(a=(1,\dots,1)\) the analogous statement is false. The counts there fall strictly below \(sn\)
once \(n\ge4\) --- by \(1\), \(2\) and \(3\) for \(\vec d=(0,1)\), \((0,1,2)\) and \((0,1,2,3)\) at
\(n=4\), and by more at \(n=5\) --- so genuine pruning occurs, and at rank two the count agrees with
the \(n+d+2\) of Conjecture \texttt{conj-ones-hrep} (Part I). Any proof of the conjecture above must therefore use
something specific to the Kozlov vector, and not merely the shape of \(\angle(a)\).

This is worth setting beside the vertex side, where the situation is reversed twice over. There the
all-ones case is the one that is proved and the Kozlov case the one that is open; here the Kozlov
polytope appears to have the simpler facet structure, and the pruning that must be understood for
\(a=(1,\dots,1)\) seems not to be needed for it at all.
\subsection{The Kozlov vector at higher rank}
\label{sec:org4369fff}
\begin{remark}
\textbf{The Kozlov vector: the same reduction, still open.} Part III of this series settles the
all-ones case by a route unavailable for a general leg vector. Passing to prefix sums works because
each coordinate enters every summand with weight one; for arbitrary \(a\) the coordinate
\(\varphi_m\) enters summand \(p\) with weight \(a_{m-d_p}\), an index relative to the window, and no
change of coordinates removes all these weights at once. This is precisely the \(a\)-dependence that
separates Theorem \texttt{thm-ones} (Part I) from Theorem \texttt{thm-kozlov-matrix} (Part I).

For the Kozlov vector both forms of the reduction nevertheless hold in every case we have examined.
Over 50 configurations --- \(a\) the Kozlov vector and \(a=(1,\dots,1)\), with \(n=3,4,5\) at rank
three and \(n=3,4\) at rank four, six shift patterns each, including ties such as
\(\vec d=(0,1,1)\) and \((0,0,1,1)\) and wide gaps such as \((0,2,4)\) and \((0,2,4,6)\) --- the
extremal tensor agreed with the conjunction of its consecutive pairwise restrictions, with no
exception. Rank four is the informative case, since there three of the six pairs are dropped rather
than one of three. Because the consecutive conditions are a subset of all pairwise conditions, which
are in turn satisfied by every index tuple in the support of \(M\), agreement of the consecutive
prediction forces agreement of all three, and the two forms were tested together.

\textbf{For the Kozlov vector this remains an observation, not a theorem.} A Minkowski sum of three or more
summands with pairwise-overlapping supports is not in general determined by its pairwise data, and
nestedness of the supports does not by itself show that the simultaneous-exposure linear program
decomposes into its pairwise subprograms. One natural certificate-level formulation can be ruled
out. By the theorem of Gordan, infeasibility of that program means some nonnegative combination of
the difference vectors vanishes; one might hope every minimal such combination to be supported on at
most two summands. It is not: minimal combinations meeting three and four summands occur for both
leg vectors, and for the Kozlov vector with \(n\ge4\) some of them use two difference vectors drawn
from a single summand. What did hold without exception, and is the statement equivalent to the
reduction, is the weaker assertion that whenever such a combination exists, one exists supported on
at most two adjacent summands. Granting the reduction, Corollary \texttt{cor-vdesc} (Part I) and Theorem
\texttt{thm-kozlov-matrix} (Part I) give a complete vertex description of \(\Koz(n;\vec d)\) for every rank.
\label{orgeb680f6}
\end{remark}
\section{Open problems for the series}
\label{sec:orgf37282b}
\begin{enumerate}
\item Give an \(H\)-representation of \(R\) for \(a=\) the Kozlov vector at rank two, matching Theorem
\texttt{thm-ones} (Part I). At that rank the vertex description is settled by Corollary \texttt{cor-vdesc} (Part I) while the
facet description is not, so this asks only for the second half of a case whose first half is
done. It is not a question about higher ranks, where even the vertex description of
\(\Koz(n;\vec d)\) is still conditional on the corresponding remark in Part III.
\item Prove Part I's conjectured \(H\)-representation for \(a=(1,\dots,1)\), facet by facet
from the vertex coordinates. The facet description of the rank-\(r\) Minkowski sum, together with
the rank-\(r\) vertex description for the all-ones simplex, is developed in Part III of this
series; the material was drafted here and moved there because the lattice-point question Part III
answers needs a precise description of the polytope. This problem is stated here because it
completes the rank-two picture of the corresponding section of Part I.
\item Prove the pairwise reduction of the corresponding remark in Part III for the Kozlov vector, which would upgrade
the vertex description of \(\Koz(n;\vec d)\) from rank two to every rank. Part III settles the
all-ones case, but by a prefix-sum argument that a general leg vector defeats, so the Kozlov case
needs a different idea. The content is that the
simultaneous-exposure linear program decomposes into its pairwise subprograms; nestedness of the
supports is suggestive but is not that statement. As recorded in the remark, the hope that every
minimal Gordan certificate meets at most two summands is false, so the target is the weaker
statement that some such certificate always meets at most two adjacent summands.
\item \textbf{A finer grading, and the lex submodules inside it.} A monomial submodule \(M\subseteq F\) makes
\(F/M\) graded by \((\mathbb Z^n)^r\), and \(H_{F/M}\) collapses all of that to a single
\(\mathbb Z\). The natural intermediate object is the \(\mathbb Z^r\)-graded Hilbert function,
which records which generator each dimension came from instead of summing over generators. Since
summing over generators is precisely the step that turns the achievable set into a sumset
(Proposition \texttt{prop-sumset} (Part I)), the multigraded refinement is exactly the invariant Theorem \texttt{thm-main} (Part I)
throws away.

At that level the question is settled, and settled cheaply. A submodule of \(F\) is
\(\mathbb Z^r\)-homogeneous exactly when it is an ids-module \(M=I_1g_1\oplus\dots\oplus I_rg_r\),
so its \(\mathbb Z^r\)-graded Hilbert function is the \(r\)-tuple of the ordinary Hilbert
functions of the \(E/I_i\), shifted. The ideals may be chosen independently, so the achievable set
is the \(r\)-fold \emph{product} of the rank-one achievable set, with no interaction constraint at all,
and the analogue of \(\Koz(n;\vec d)\) is a product rather than a Minkowski sum. Applying the
exterior Kruskal--Katona theorem in each component separately shows that every
\(\mathbb Z^r\)-graded Hilbert function is realized by a tuple of lex ideals.

What is \emph{not} settled is the shape of a distinguished subset. A tuple of lex ideals is in general
not a lex submodule in the sense of Amata and Crupi, whose definition also imposes the nesting
condition \((e_1,\dots,e_n)^{\rho_i+d_i-d_{i-1}}\subseteq I_{i-1}\), where \(\rho_i\) is the
initial degree of \(I_i\); for \(n=3\) and \(\vec d=(0,0)\) only \(43\) of the \(100\) pairs of
lex ideals satisfy it. \textbf{Determine the set of \(\mathbb Z^r\)-graded Hilbert functions of
Amata--Crupi lex submodules inside that product, and in particular its convex hull.} The question
is natural because lex submodules are in bijection with the achievable \(\mathbb Z\)-graded
Hilbert functions, so this set is a canonical section of the collapsing map from the product to
the sumset --- a choice, for each achievable \(\mathbb Z\)-graded Hilbert function, of one
preferred way to split it among the generators. Its convex hull would be the multigraded analogue
of the Kozlov polytope itself. The nesting condition is a chain condition on initial degrees, so a
staircase or order-polytope answer is the thing to look for. Whether this belongs here or in a
later part of the series is undecided.
\item \textbf{Sköldberg's equivalence at the module level.} His sign-twisting equivalence between modules over
the exterior face ring and over the artinified Stanley--Reisner ring extends to an isomorphism
\(\mathrm{AIM}(\vec\Delta)\cong\IM(\vec\Delta)\) for indicator modules, which should make Betti
numbers, not just Hilbert functions, transportable between the two settings.
\item \textbf{Lifting Gröbner bases.} Nilsson and Snellman (2003) lift a Gröbner basis of an ideal
\(I\subseteq E\) to one of its preimage in the free associative algebra, via a section map plus a
"squeezing" correction. Whether the same construction works for submodules of \(F\) under a
position-over-term module order is open, and the obstacle is concrete: squeezing is defined purely
by a monomial's index gap, i.e. entirely in terms of term parts, while a position-over-term order
never compares term parts across components.
\end{enumerate}

Beyond these, the question this paper deliberately does not treat --- which \emph{lattice points} of
\(\Koz(n;d_1,\dots,d_r)\) are achievable --- is the subject of the companion paper, where it is
answered by a repaired form of Amata and Crupi's generalized Kruskal--Katona theorem. The two
questions are related by Example \texttt{ex-worked} (Part I): the hull is a coarse object, and an achievable Hilbert
function can sit well inside it.
\end{document}
